% =============================================================================
% From Speakers to Dialects
% Cristina Gutierrez Alcala - Speech Signal Processing, 2026
%
% ACL style. Change the \usepackage line below to match the style file you
% already have in Overleaf (acl.sty, ACL2023.sty, emnlp2023.sty, ...).
% Upload the seven PNGs from output/figures/ into the Overleaf project root.
% =============================================================================

\documentclass[11pt]{article}

\usepackage[]{acl}

\usepackage{times}
\usepackage{latexsym}
\usepackage[T1]{fontenc}
\usepackage[utf8]{inputenc}
\usepackage{microtype}
\usepackage{graphicx}
\usepackage{booktabs}
\usepackage{makecell}
\usepackage{amsmath}
\usepackage{inconsolata}

\title{From Speakers to Dialects:\\An Acoustic Study of Northern and Southern
Peninsular Spanish Using Common Voice}

\author{Cristina Guti\'errez Alcal\'a \\
  Speech Signal Processing \\
  2026}

\begin{document}
\maketitle

\begin{abstract}
This study compares Northern and Southern Peninsular Spanish acoustically using
dialectally labeled recordings from Mozilla Common Voice, with forced alignment
for phonemic measurement and statistical inference at the speaker level. Across
957 speakers and 89,958 vowel tokens, no difference was detected between the two
varieties in the vowel space: no value exceeded $|d| = 0.17$ in a design powered
for $d \geq 0.24$. The fact that the result is attributable to the dialects and
not to the method is corroborated by an external Mexican group, which does show a
clear separation from both Peninsular varieties with effects 2.5 times greater.
The variation is localized in the consonant: Southern speech weakens /s/ in coda
($d = +0.63$), with /s/ at the onset serving as the null control. The study also
documents how an earlier version of this same analysis yielded effects exceeding
one standard deviation that turned out to stem from two individual speakers, and
what is required of a volunteer corpus to avoid reproducing that error.
\end{abstract}

% =============================================================================
\section{Introduction}

The fact that a speaker from Bilbao and one from Seville sound different is an
intuition shared by all Spanish speakers and one that dialectology has described
for over a century. Translating that intuition into a reproducible acoustic
measure is a problem that remains largely unresolved, and it is the one addressed
in this paper.

The first two formants of a vowel, F1 and F2, are the resonance frequencies of
the vocal tract that directly encode its height and frontness, and these
constitute the standard representation of a variety's vowel space. The initial
question was whether the vowel space of Northern peninsular Spanish differs
measurably from that of Southern peninsular Spanish, and whether that difference
is sufficient to classify an unknown speaker. Mozilla Common Voice provided an
unusual resource for answering these questions: thousands of recordings with
dialect labels declared by the speakers themselves, openly available and on a
scale that laboratory phonetics rarely achieves.

An initial version of the analysis answered ``yes'', and emphatically so:
differences exceeding one and a half standard deviations in F1 and in the
bandwidth of F2, with $p$-values below 0.001. A subsequent audit showed that
those effects stemmed from two specific speakers, who together accounted for
54\% and 80\% of their respective dialect cells. The study did not compare two
varieties; it compared two people.

This report presents an analysis reconstructed from that finding incorporating it
as the subject of study rather than as an error to be concealed. The study
establishes four findings. First, it documents the mechanism by which clip-based
sampling generated large, spurious effects, and shows that the magnitude obtained
is a direct function of the extent to which an individual is allowed to dominate
the sample (Section~\ref{sec:pseudo}). Second, using a balanced design with 957
speakers, it finds no difference whatsoever between the two peninsular varieties
in the vowel space, in a contrast with sufficient power to rule out anything
other than a small effect (Section~\ref{sec:vowelspace}). Third, it verifies that
this absence is not due to the coarseness of the method by incorporating Mexican
Spanish as an outgroup, which does separate from both varieties using the same
procedure (Section~\ref{sec:outgroup}). And fourth, it finds a substantial effect
in the consonant system: the weakening of /s/ in coda in Southern speech
(Section~\ref{sec:coda}).

The common thread linking the four is a matter of design rather than statistics,
and is stated in Section~\ref{sec:discussion}: in a corpus built from unequal
contributions by volunteers, the contrasts identified within individual speakers
are robust, while the overall comparisons between groups are fragile.

% =============================================================================
\section{State of the art}

\subsection{Dialectal variations in Peninsular Spanish}
\label{sec:dialectal}

Spanish has a remarkably stable system of five vowels: /a, e, i, o, u/. It lacks
distinctions based on vowel length and tension; its unstressed vowels are not
reduced as they are in English or Portuguese; and its inventory has hardly
changed since the Romance language system became fixed. This stability has a
direct implication for this study: the vowel system offers little scope for
dialectal differentiation, and variation in Spanish is overwhelmingly consonantal
\citep{hualde2005,martinez2007}. The reference formant values for peninsular
Spanish were established by \citet{quilis1981} based on laboratory speech and
remain the standard basis for comparison.

The best-documented feature of Southern speech is the weakening of /s/ in
implosive position, which is realized as aspiration, assimilation to the
following segment, or elision, in contrast to the characteristic retention found
in the peninsular North
\citep{alvar1996,narbona1998,gerfen2002,ruch2014}. Associated with this
weakening, a classic descriptive account holds that the preceding vowel
compensates for the loss by opening and lengthening, to the point of positing an
expanded vowel system based on degree of openness for Eastern Andalusian
\citep{alonso1950}. This assertion, based on auditory observation and small
samples, can be tested directly with current methods, and
Section~\ref{sec:precoda} does so.

\subsection{Formant Normalization}
\label{sec:norm}

Formant frequencies depend on the length of the vocal tract as much as on
articulatory configuration, so that two speakers producing the same vowel will
systematically differ in hertz if they differ in anatomy. Comparing varieties in
raw hertz values, when the data sets are not balanced by sex, is to confuse
dialect with physiology. Sociophonetics solves this problem by normalizing data
per speaker, and among the available procedures, the intra-speaker z-score
proposed by \citet{lobanov1971} performs best in systematic comparisons: it
eliminates anatomical variation while preserving sociolinguistic differences
\citep{adank2004}. The Bark scale \citep{traunmuller1990}, which is perceptually
based, is used here as a complementary measure.

\subsection{Common Voice as a phonetic corpus}
\label{sec:corpus}

Common Voice is a multilingual corpus of read speech built from contributions by
volunteers and distributed openly \citep{ardila2020}. For dialectology, it offers
something rare: variety labels declared by the speakers themselves, on a scale of
thousands of speakers. It also presents two challenges. The first is that these
labels are neither verifiable nor granular. The second, more significant
challenge is that the distribution of contributions is extraordinarily uneven,
because nothing in the project's design encourages equitable distributions: a few
enthusiastic volunteers can each contribute hundreds of recordings. Sampling that
selects clips without grouping them by speaker therefore produces observations
that are not independent, with the consequent risk of pseudoreplication
\citep{hurlbert1984,winter2019}. The availability of high-quality forced aligners
for Spanish \citep{mcauliffe2017} is what allows us to move from clip-level
measures to phoneme-level measures, and thereby to examine the intra-speaker
contrasts that the corpus does indeed robustly support.

% =============================================================================
\section{Corpus and sampling}

\subsection{Common Voice 26.0 and dialect tags}

All the data used come from Mozilla Common Voice, Spanish 26.0 version (12 of
June of 2026), a corpus of read speech built from voluntary submissions and
distributed under an open licence \citep{ardila2020}. The version used contains
439,225 validated clips. Each clip has an orthographic transcription associated
with it and a speaker identification number (\texttt{client\_id}), in addition to
autodeclared metadata such as gender, age and accent.

Getting the corpus onto a machine where I could work with it shaped the pipeline
more than I expected. The Spanish release is a single archive of roughly 48 GB.
My first attempt was to download it, unpack it and work from the audio. That cost
two days of failed transfers and still left me without room for the extracted
files. The fix turned into the design. Nothing in the sampling stage needs the
audio, because the selection is made entirely from \texttt{validated.tsv}, which
sits at the front of the archive and can be pulled out with
\texttt{tar -{}-fast-read} in about a second. Only the clips the sampler selects
are then extracted, which in the final study meant unpacking 5,577 files instead
of 439,225.

The accent tag is what makes this corpus usable in a dialectal study, since the
speakers are responsible for tagging the dialectal variation that they consider
they speak. From the available Spanish tags, I selected three: ``peninsular
septentrional'' also referred to as North, ``peninsular meridional'' or South and
Mexican. The latter being incorporated as an external group due to the reasons
explained in Section~\ref{sec:outgroup}. The speakers are also able to mark one
or more accents, and a non-negligible proportion do. I decided to discard those
clips instead of manually assigning one of the tags, since any tie-breaker
criterion would introduce a biased opinion in the grouping variable. After the
filter was applied, I had 236,832 clips with exactly one of the aforementioned
tags.

The nature of these tags must also be mentioned. They are autodeclarations and
not dialectological classifications. There is no independent verification of the
provenance of the speaker, or further fine-grained information than the region.
This limitation will be touched on in Section~\ref{sec:limitations}.

\subsection{Speaker-level sampling and per-speaker cap}

The most important design choice of this project was also the simplest: the
sampling was done directly on the speaker and not on the clip. An older version
of this analysis selected clips randomly, identifying them by file path, without
looking at who had produced each clip. The result was a sample dominated by a
handful of individuals. Section~\ref{sec:pseudo} is devoted entirely to the
consequences.

The current sampling groups first by \texttt{client\_id}, arranges the clips of
each speaker by a reproducible criterion and applies a cap of eight clips per
speaker. The declared design objective is to first maximize the number of
speakers and then the number of clips per speaker, because that same number of
speakers (and not number of tokens) is what narrows down inference. The whole
sampling stage is seeded, and is registered in a versioned manifest, so that the
sample is reconstructible clip by clip.

The cap of eight is a compromise between two requirements that pull in opposite
directions. Lobanov normalization needs enough vowels from each speaker to
estimate that speaker's own mean and spread, which argues for more clips per
person. Independence argues for fewer. Eight clips of roughly five seconds yield
on the order of a hundred vowel tokens per speaker, comfortably enough for the
normalization, while holding any single individual's share of a dialect cell
below one percent. The sampling is implemented with pandas \citep{mckinney2010}
over NumPy arrays \citep{harris2020}.

One consequence of prioritizing speakers is that the design is bounded by its
smallest group. Southern Peninsular Spanish has 188 speakers in release 26.0, and
that number constrains everything the study can claim, so Northern and Southern
speakers were both taken in full. Mexican Spanish has over 1,500 available, far
more than the comparison needs, and was capped at 300.

The effect of the cap is measured with the Gini coefficient of the distribution
of clips per speaker, where 0 indicates a perfectly even distribution and 1 a
total concentration. In the preliminary sample the Gini coefficient was 0.76 in
North and 0.85 in South. In the current sample it is 0.26 in North and 0.18 in
South. This sole change is what made the rest of the analysis interpretable.

\subsection{Final design of the sample}
\label{sec:sample}

The resulting sample has 5,577 clips of 957 different speakers, and a total of
484.5 recorded minutes. Table~\ref{tab:sample} compiles its composition.

\begin{table}[t]
\centering
\small
\setlength{\tabcolsep}{4pt}
\begin{tabular}{lrrrrr}
\toprule
Variety & Clips & Speakers & \makecell{With\\tokens} & \makecell{Norma-\\lizable} & Tokens \\
\midrule
North  & 2,639 & 469 & 468 & 443 & 42,805 \\
South  & 1,199 & 188 & 188 & 181 & 19,549 \\
Mexico & 1,739 & 300 & 299 & 282 & 27,604 \\
\midrule
Total  & 5,577 & 957 & 955 & 906 & 89,958 \\
\bottomrule
\end{tabular}
\caption{Composition of the sample in each variety.}
\label{tab:sample}
\end{table}

The imbalance between varieties is dictated by the corpus itself. The Southern
part of the peninsula is represented by fewer than half as many speakers as the
North. I do not correct for this by undersampling the North, because discarding
valid data solely to balance the cells would reduce statistical power without
gaining anything. The comparisons are calculated using tests that do not assume
equal sample sizes.

The imbalance by sex, however, does matter. All three varieties are dominated by
male voices: 27.9\% of female speakers in the North, 20.4\% in the South, and
25.2\% in Mexico, with the rest being male or undeclared. Since the length of the
vocal tract directly influences formant frequencies, a comparison between these
groups in raw hertz does not distinguish dialect from anatomy. This is why the
normalization described in Section~\ref{sec:qc} is not an optional refinement but
a design requirement.

Figure~\ref{fig:dataflow} summarizes the complete path, from the 439,225
validated clips in the corpus to the tokens finally analyzed. The alignment and
quality-control stages that produce those numbers are described in
Sections~\ref{sec:align} and~\ref{sec:qc}.

\begin{figure*}[t]
\centering
\includegraphics[width=\textwidth]{fig_data_flow.png}
\caption{From corpus to analyzed tokens. The left panel shows the reduction from
the validated Common Voice clips to the final sample after filtering by unique
label and applying the limit of eight clips per speaker. The right panel shows
the quality control funnel applied to the extracted vowel tokens.}
\label{fig:dataflow}
\end{figure*}

The sample of 903 clips on which the preliminary analysis was based remains
unchanged and is reused in Section~\ref{sec:pseudo}, where it serves not as a
data source but as an object of study.

% =============================================================================
\section{Methodology}

\subsection{Audio preprocessing and forced alignment}
\label{sec:align}

Each clip was converted from MP3 to mono WAV at 16 kHz, which is the sample
frequency expected by the aligner, with peak scaling to avoid clipping. Next to
each audio file, a text transcription file was written from what the speaker read
aloud, taken directly from the corpus metadata, and the ensemble was organized in
per speaker directories so that the aligner could apply its speaker adaptation.

Forced alignment was not part of the original plan. It became necessary after two
attempts to recover vowel identity by cheaper means, and both failures explain
why the pipeline looks the way it does. The first classified each measurement by
its position in the F1/F2 plane, using the reference ranges for Spanish vowels in
\citet{quilis1981}. Nearly every token it labelled /a/ came from the Northern
group and nearly every /i/ from the Southern one. Those ranges are calibrated on
Northern speech, so a displaced Southern token lands in a neighbouring category.
The procedure is also circular: labelling vowels by their formants and
then comparing formants across dialects builds the conclusion into the labels.
The second attempt inferred the vowel from the transcription, which works only
for clips containing a single vowel type. Of 903 clips, 14 resolved.

That left aligning the audio to its transcription properly. I installed the
Montreal Forced Aligner \citep{mcauliffe2017}, built on Kaldi \citep{povey2011},
together with the pretrained \texttt{spanish\_mfa} acoustic model and dictionary.
The aligner takes the audio and its orthographic transcription and returns a
TextGrid with the time boundaries of each word and each phone. Of the 5,577
clips, 5,572 aligned correctly (99.9\%); the remaining five failed because of
discrepancies between the transcription and the signal.

With alignment, each clip yields all of its vowel nuclei with phonemic labels,
and with them four pieces of information that were unavailable before: vowel
identity, duration, stress position, and the following segment. The last of these
is what made the coda /s/ analysis of Section~\ref{sec:coda} possible. The
semivowels /j/ and /w/ are deliberately excluded, since they are not syllable
nuclei and their formant targets are transitional.

\subsection{Extraction of acoustic parameters}

The acoustic measures were taken using Praat \citep{boersma2021} through its
Python interface Parselmouth \citep{jadoul2018}, using linear predictive analysis
with Burg's method with a maximum of five formants.

The frequency ceiling is adapted to the speaker's declared sex: 5,000 Hz for male
voices, 5,500 Hz for female voices, and 5,250 Hz when the corpus records no
value. The previous analysis applied a singular ceiling of 5,000 Hz to all of the
speakers, which systematically biased the estimation of female formants. This was
not a trivial problem, since the dialect cells are not balanced by sex.

From each vowel, F1, F2, B1 and B2 were registered in three different points:
25\%, 50\%, and 75\% of its duration. The middle point is the conventional
measure; the other two allow to quantify the formant movement inside a vowel,
relevant to descriptions of diphthongization in Southern speech. To them, I added
segment duration, fundamental frequency, the harmonics-to-noise ratio (HNR), the
clip's signal-to-noise ratio, lexical stress and the following consonantal
context. The vowels with a duration less than 30 ms were discarded from the
start, since they were already dominated by the transitions towards the adjacent
consonants and lacked a measurable stationary state. The extraction was
parallelized on various processes.

\subsection{Quality control and normalization}
\label{sec:qc}

Quality control first applies deliberately generous range filters: F1 between 150
and 1,100 Hz, F2 between 500 and 3,200 Hz, and bandwidth between 10 and 900 Hz
(B1) and 10 and 1,000 Hz (B2). These limits exist to remove tracking failures (an
F1 of 3 kHz is an octave error, not a vowel), not to impose a phonetic
expectation. Tightening them further would reintroduce the circularity that
invalidated the vowel tagging of the preliminary analysis. Also, F2 must be
higher than F1 and outliers beyond three standard deviations within each speaker
and each vowel were eliminated, a filter that is only possible when the identity
of the segment is known.

\begin{table}[t]
\centering
\small
\begin{tabular}{lrr}
\toprule
Phase & Tokens & \% retained \\
\midrule
Extracted                        & 100,053 & 100.0 \\
Measurable at midpoint           & 100,052 & 100.0 \\
F1/F2 in range                   &  99,104 &  99.1 \\
B1/B2 in range                   &  91,198 &  91.1 \\
F2 $>$ F1                        &  91,198 &  91.1 \\
Within-speaker outliers removed  &  89,958 &  89.9 \\
\bottomrule
\end{tabular}
\caption{Quality control funnel on the extracted vowel tokens.}
\label{tab:qc}
\end{table}

The main loss is concentrated in the bandwidth filters, which discards 8\% of
tokens. The bandwidths estimated by linear prediction were notably unstable and
sensitive to the recording conditions, a fact that is relevant in
Section~\ref{sec:pseudo}, where one of the preliminary analysis effects was
precisely an effect of B2.

On the surviving tokens, three representations were calculated. The first is raw
hertz, which is retained for transparency but does not support any comparisons
between groups, for the reasons explained in Section~\ref{sec:sample}. The second
is the Bark scale \citep{traunmuller1990}, which approximates the ear's frequency
resolution. The third, and the one used in all contrasts, is Lobanov
normalization \citep{lobanov1971}: the z-score of each formant computed within
each speaker across all of their vowels. It is the best-performing procedure in
systematic comparisons of normalization methods \citep{adank2004}, for the
reasons set out in Section~\ref{sec:norm}. It requires that the speaker has
produced all five vowel categories, a condition met by 906 of the 955 speakers
with tokens. The remaining 49 are excluded from the normalized analyses.

\subsection{Statistical scope}

The principle that governs all inference is that the unit of analysis is the
speaker. Tokens from the same speaker are not independent observations. They all
share the same anatomy, microphone, recording environment, and reading style.
Treating them as independent inflates the degrees of freedom and produces
significance regardless of whether or not a dialectal effect exists, as
demonstrated empirically in Section~\ref{sec:pseudo}.

In practice, this means that the data are aggregated into speaker $\times$ vowel
cells and, from there, into a single measure per speaker, prior to any
comparisons. Comparisons between varieties are performed using Welch's $t$-test,
which does not assume equal variances or sample sizes, along with Cohen's $d$
\citep{cohen1988} and confidence intervals obtained via bootstrapping with 2,000
resamples of speakers, not tokens. Welch's test and the bootstrap are computed
with SciPy \citep{virtanen2020}, and the mixed models are fitted with statsmodels
\citep{seabold2010}. Within each family of contrasts, $p$-values are corrected
for multiple comparisons using Holm's procedure \citep{holm1979}. Mixed models
with a random intercept per speaker are fitted to the means per speaker and
vowel. The significance level is 0.05, and every stochastic step, like sampling,
bootstrapping, $k$-means initialization, and cross-validation partitions, is
seeded with a fixed seed.

A methodological caveat is worth noting, because it is an easy pitfall to fall
into. Fitting a token-level mixed model to values already normalized by Lobanov
yields an ICC close to zero and $p$-values less than 0.001. Lobanov sets each
speaker's mean to 0 and their standard deviation to 1, so that there is no
remaining variance between speakers for a random intercept to account for, and
the model reverts to treating the tokens as independent. A random intercept
protects against speakers differing in level, not against one speaker shaping a
within-speaker pattern. Aggregating into speaker $\times$ vowel cells before
modeling is what solves the problem.

Finally, the design is accompanied by a power analysis computed from the realized
speaker counts, which sets limits on what the study can honestly claim. The
resulting thresholds are reported alongside the results in
Section~\ref{sec:vowelspace}.

\subsection{Validation of the extractor}

To verify that the extractor measures what it claims to measure, vowels with
known formants are synthesized using an F1--F5 resonator bank on a 120 Hz glottal
source, and the recovered value is compared with the actual value. Measurement is
accurate for open vowels (the error in F1 for /a/ is +0.2\%) but degrades for
closed vowels, with an error of +13.5\% for /i/ and +33.7\% for /u/. The cause is
well known: in closed vowels, F1 approaches F0, so that individual harmonics
dominate the fitted envelope, and in /u/, the F1 and F2 poles interfere with each
other.

The bias is identical in both dialects, so it cannot generate a spurious
difference between them, but it makes /u/ the least reliable measure in the
study, and it is treated as such in the discussion. The test is part of a battery
of 24 automated tests that also verify the handling of coda context, the
identification of the stressed syllable, and the normalization functions.

The pipeline runs as a single seeded entry point; figures are produced with
Matplotlib \citep{hunter2007}, and TextGrids are parsed with praatio. A full run
takes about fifteen minutes end to end.

% =============================================================================
\section{Pseudoreplication}
\label{sec:pseudo}

\subsection{Naive analysis}

The first version of this study worked on 440 tokens (244 North and 196 South),
obtained from taking a measurement per clip in the maximum global intensity peak.
Applying Welch $t$-test to those tokens, three out of four comparisons were
significant, and two of them with effects of a magnitude that is rarely observed
in dialectal phonetics. F1 had a $d = +1.34$ and B2 $d = -1.43$, both with
$p < 0.001$. Literally interpreted, these numbers described two varieties whose
vowels were separated by more than one and a half standard deviations.

The result was, in appearance, excellent. It was also false, and the way in which
it was false is more instructive than any effect it might have found.

\subsection{Audit per speaker}

The 440 tokens do not come from 440 different speakers. When the measurements are
cross-referenced with the speaker identifiers of the corpus, the sample revealed
54 Northern speakers and 26 Southern speakers, with an extraordinarily unequal
distribution. A single speaker contributed 132 tokens, which was 54\% of the
septentrional cell, and another one made up 157, 80\% of the meridional cell. The
three most prolific speakers of each cell concentrated 64\% and 86\% of the data,
respectively.

The consequences can be directly seen by eliminating speakers one by one. If the
dominant meridional speaker is discarded, $d$(F1) falls from 1.34 to 0.23 and
$d$(B2) from $-1.43$ to $-0.45$. If the septentrional speaker is discarded, they
fall to 1.04 and $-1.04$. If both are removed (289 of the 440 tokens, leaving
151), the effects are reduced to $d = +0.01$ and $d = -0.09$, translated to
almost nothing. The elimination of any other speaker leaves the numbers
practically intact.

What the preliminary analysis measured, thus, was not the difference between
septentrional and meridional Spanish, but the difference between two concrete
people. This is the problem that \citet{hurlbert1984} called pseudoreplication.
Pseudoreplication means treating as independent replicas observations which share
the experimental unit. The $t$-test assumes independence; 440 clips from 80
speakers do not satisfy it, and when one of these makes up 80\% of a cell, the
effective sample size of that cell gets closer to one than to two hundred.

\subsection{Size of the effect}

The dependence can be directly quantified by recalculating the original analysis
under increasing token caps per speaker. With a cap of one token per speaker (80
strictly independent observations) $d$(F1) is $+0.11$. With a cap of three,
$+0.11$. With a cap of eight, $+0.30$. With a cap of twenty, $+0.54$; with a cap
of fifty, $+0.90$; with a cap of one hundred, $+1.15$; and with no cap at all,
$+1.34$. The relationship is monotonic, and it is replicated in B2 with the
opposite sign.

In other words: the published effect size is a direct function of how much an
individual is allowed to dominate the sample. There is no ``true'' value that the
cap distorts; the cap is the sole determinant of the obtained value.

\subsection{Naive against corrected}

Table~\ref{tab:naive} contrasts the four original comparisons with two analyses
that respect the data structure: the aggregation of one measure per speaker and a
mixed model with a random intercept per speaker \citep{barr2013,winter2019}.

\begin{table}[t]
\centering
\small
\setlength{\tabcolsep}{4pt}
\begin{tabular}{lrrrrrr}
\toprule
 & \makecell{naive\\$d$} & \makecell{$p$\\(Holm)} &
 \makecell{per spk\\$d$} & $p$ & \makecell{mixed\\$p$} & ICC \\
\midrule
F1 & $+1.335$ & $<$0.001 & $+0.040$ & 1.000 & 1.000 & 0.54 \\
F2 & $-0.337$ & 0.001    & $+0.333$ & 0.835 & 1.000 & 0.17 \\
B1 & $-0.166$ & 0.100    & $+0.195$ & 1.000 & 1.000 & 0.05 \\
B2 & $-1.433$ & $<$0.001 & $-0.172$ & 1.000 & 1.000 & 0.28 \\
\bottomrule
\end{tabular}
\caption{Naive analysis per token against a speaker level analysis, on 440
tokens.}
\label{tab:naive}
\end{table}

No effect survives. Note also that F2 not only loses its significance but also
changes sign when moving to the speaker level, which rules out interpreting it as
an attenuated real effect.

The ICC for F1 is 0.54: more than half of the variance in F1 is explained by who
is speaking, not by which dialect they speak. From this value, the design effect
can be calculated, indicating that the 440 pooled tokens carry the information of
approximately 128 independent observations. The bootstrap-constructed confidence
intervals show the same result from another angle. When resampling tokens,
$d$(F1) has an interval of $[+1.11, +1.60]$, which comfortably excludes zero.
When resampling speakers, the interval is $[-0.20, +1.83]$ and includes it. For
B2, $[-1.67, -1.22]$ versus $[-1.89, +0.19]$. All four speaker-resampled
intervals contain zero.

\begin{figure*}[t]
\centering
\includegraphics[width=0.75\textwidth]{fig_pseudoreplication.png}
\caption{Anatomy of a sampling artifact, based on the 440 tokens from the
preliminary analysis. A: A few speakers contribute most of the data for each
dialect. B: The effects fade as the cap on tokens per speaker prevents any single
individual from dominating. C: Removing a single prolific speaker eliminated the
effects. D: The confidence intervals obtained by resampling speakers all include
zero, unlike those obtained by resampling tokens.}
\label{fig:pseudo}
\end{figure*}

\subsection{Why the balanced result helps}

When applied to the sample described in Section~\ref{sec:sample}, the same audit
procedure finds nothing to audit. The most prolific southern speaker now accounts
for 1\% of their cell, and their removal shifts the largest per-vowel effect by
0.00, compared to the shift from 1.6 to 0.0 produced by the same operation in the
previous sample. The pooled analysis and the speaker-level analysis coincide,
simply because there is nothing left on which they might disagree.

This convergence is the argument in favor of the per-speaker cap, and it is also
the reason why the results in Section~\ref{sec:results} can be interpreted with
confidence: not because the statistical method is more sophisticated, but because
the sampling design no longer allows the question to answer itself.

% =============================================================================
\section{Results}
\label{sec:results}

\subsection{Vowel space on speaker level}
\label{sec:vowelspace}

The main comparison of this project was made on a per-speaker and per-measure
average, between 456 and 461 Northern speakers and between 178 and 186 Southern
speakers, according to how many produced the vowel at hand.
Table~\ref{tab:speakerlevel} compiles the fifteen evaluated measures: the
normalized positions of the five vowels in F1 and F2, the average duration, the
bandwidths, the harmonics-to-noise ratio and the fundamental frequency.

\begin{table}[t]
\centering
\small
\setlength{\tabcolsep}{4pt}
\begin{tabular}{lrrrrr}
\toprule
Measurement & North & South & $d$ & $p$ & \makecell{$p$\\(Holm)} \\
\midrule
F1 /a/       & 1.09   & 1.07   & $+0.12$ & 0.166 & 1.000 \\
F1 /e/       & $-0.31$ & $-0.28$ & $-0.17$ & 0.043 & 0.638 \\
F1 /i/       & $-1.22$ & $-1.21$ & $-0.03$ & 0.720 & 1.000 \\
F1 /o/       & $-0.09$ & $-0.08$ & $-0.04$ & 0.661 & 1.000 \\
F1 /u/       & $-0.84$ & $-0.84$ & $+0.01$ & 0.893 & 1.000 \\
F2 /a/       & $-0.25$ & $-0.27$ & $+0.14$ & 0.085 & 1.000 \\
F2 /e/       & 0.59   & 0.58   & $+0.02$ & 0.783 & 1.000 \\
F2 /i/       & 1.50   & 1.52   & $-0.06$ & 0.440 & 1.000 \\
F2 /o/       & $-0.96$ & $-0.95$ & $-0.04$ & 0.684 & 1.000 \\
F2 /u/       & $-1.10$ & $-1.16$ & $+0.11$ & 0.135 & 1.000 \\
Duration (s) & 0.077  & 0.077  & $-0.00$ & 0.995 & 1.000 \\
B1 (Hz)      & 149.2  & 146.1  & $+0.05$ & 0.528 & 1.000 \\
B2 (Hz)      & 218.0  & 216.8  & $+0.02$ & 0.822 & 1.000 \\
HNR (dB)     & 9.97   & 9.81   & $+0.04$ & 0.616 & 1.000 \\
F0 (Hz)      & 151.3  & 149.4  & $+0.05$ & 0.580 & 1.000 \\
\bottomrule
\end{tabular}
\caption{Speaker level comparison between North and South. F1 and F2 values are
normalized by Lobanov (intra-speaker z-scores); the bandwidths, the duration, the
HNR and F0 are expressed using their original units. A positive value of $d$
indicates a higher value in the North.}
\label{tab:speakerlevel}
\end{table}

None of these comparisons reaches significance after the correction by multiple
comparisons. The major observed effect in any point on this table is
$|d| = 0.17$, corresponding to an F1 of /e/, whose $p$ value without correcting
(0.043) stays at 0.638 after applying Holm on the fifteen contrast family. The
bandwidths, where the preliminary analysis located its second strongest effect
($d = -1.43$), now show $d = +0.02$.

The same conclusion is obtained with global vowel space measurements instead of
using individual positionings. The area of the convex hull defined by the five
normalized vowels is 3.05 in the North and of 2.95 in the South, with $d = +0.12$
and a confidence interval at 95\% of $[-0.05, +0.28]$, $p = 0.167$. Mixed models
fitted to the speaker $\times$ vowel means show no significant main effect of
dialect.

\begin{figure*}[t]
\centering
\includegraphics[width=\textwidth]{fig_vowel_space.png}
\caption{Vowel space of the Northern and Southern parts of the Iberian Peninsula.
On the left, in raw hertz; on the right, after Lobanov normalization. Each point
represents the average for a single speaker, and the ellipses cover 95\% of the
variation among speakers. The reduction in dispersion between the two panels
corresponds to the anatomical variation that normalization eliminates; what
remains is nearly a complete overlap between varieties.}
\label{fig:vowelspace}
\end{figure*}

What makes this result informative is not the absence of significance, but the
power of the design. With 461 Northern speakers and 186 Southern speakers, the
comparison would detect, with 80\% probability, any effect of magnitude
$d \geq 0.24$, and with 90\% probability, any effect of $d \geq 0.28$. The sample
from the preliminary analysis, by contrast, could only detect effects of
$d \geq 0.63$, that is, large effects only.

The conclusion drawn from these data must be formulated with precision. It is not
that the vowel spaces of both varieties are identical: differences below
$d \approx 0.24$ remain entirely possible, and this study would not rule them
out. What can be stated, however, is that there is no difference in the vowel
space between the North and South of the Peninsula large enough to be detected by
a design with more than six hundred speakers. Any effect that exists is, at most,
small.

The obvious objection remains: that the method is simply incapable of
distinguishing dialects. The following section addresses this objection.

\subsection{Calibration with an external group}
\label{sec:outgroup}

The result in Section~\ref{sec:vowelspace} allows for two incompatible
interpretations. One is that the two peninsular varieties do not differ
appreciably in their vowel systems. The other is that the procedure is too crude
to detect any dialectal differences, and would have produced a null result
regardless of the comparison made. Without additional evidence, both explanations
account for the data equally well.

The way to distinguish between them is to apply exactly the same procedure to a
contrast whose existence is not in dispute. To that end, 300 Mexican speakers
were included in the design and processed using the same sampling, the same
alignment, the same feature extraction, the same normalization, and the same
statistical tests. If the instrument distinguishes between Peninsular and Mexican
Spanish but not between Northern and Southern Spanish, the absence of an effect
in Section~\ref{sec:vowelspace} is a statement about the dialects, not about the
method.

The three pairwise contrasts, based on the twelve vowel measures evaluated at the
speaker level, show the following results. North versus South gives an average
$|d|$ of 0.064 and a maximum of 0.172, with 0 of the 12 comparisons significant.
North versus Mexico gives 0.152 and 0.303, with 4 of 12. South versus Mexico
gives 0.173 and 0.411, with 2 of 12.

Transatlantic contrasts produce effects that are, on average, 2.5 times greater
than those of the Peninsular contrast, and reach statistical significance after
correction in six of the twenty-four comparisons, whereas the North-South
contrast does not reach significance in any of the twelve. The largest
differences are concentrated, as might be expected, in F2 of the peripheral
vowels: F2 of /a/ between the South and Mexico ($d = -0.41$, $p < 0.001$), F2 of
/i/ between the South and Mexico ($d = +0.40$, $p < 0.001$), and F2 of /i/
between the North and Mexico ($d = +0.30$, $p = 0.004$). Vowel duration also
distinguishes the North from Mexico ($d = -0.28$, $p = 0.032$), without
distinguishing the North from the South.

\begin{figure*}[t]
\centering
\includegraphics[width=\textwidth]{fig_outgroup.png}
\caption{Calibration with an external group. The three varieties overlaid in the
normalized vowel space and the average magnitude of the effect for each pairwise
contrast. The peninsular contrast is consistently the smallest of the three.}
\label{fig:outgroup}
\end{figure*}

The conclusion is that the method does resolve dialectal differences when they
exist, using the same sample, the same measures, and the same decision threshold.
The null result in Section~\ref{sec:vowelspace} is not an artifact of a blunt
instrument.

It is worth saying here, rather than in the discussion, that these effects are
small in absolute terms. A mean $|d|$ of 0.16 for the transatlantic contrast is
not what the descriptive literature would lead you to expect from two varieties
separated by an ocean and five centuries. I read that as a problem of scale
rather than a failure of the comparison. The five-vowel system described in
Section~\ref{sec:dialectal} gives dialectal drift very little to work with, and
the outgroup is what makes this visible: even across the Atlantic, these vowels
barely move. Seen that way, the null result of Section~\ref{sec:vowelspace} is
close to what Spanish phonology would predict. This immediately suggests where
variation should be sought instead of in the vowels. The following section
explores this.

\subsection{/s/ in coda}
\label{sec:coda}

If dialectal variation in Spanish lies primarily in the consonants, the obvious
candidate for contrast within the Iberian Peninsula is the behaviour of /s/ in
syllable-final position. Descriptions of the Southern speech consistently point
to its weakening (aspiration, assimilation to the following segment, or elision)
in contrast to the full realization retained by Northern speech
\citep{narbona1998,gerfen2002,torreira2007}.

This hypothesis has a feature that makes it particularly suitable in light of
what we saw in Section~\ref{sec:pseudo}: it is contrasted within each speaker.
The relevant comparison is not between the Northern /s/ and the Southern /s/, but
between the /s/ in coda and the /s/ at the onset of the same speaker. Each
speaker thus acts as their own control, and an imbalance in the number of clips
per speaker cannot produce the effect, because it affects both terms of the
comparison equally. This is exactly the type of contrast that the preliminary
analysis could not address, due to the lack of phonemic segmentation.

The alignment yields 21,639 /s/ segments from 890 speakers, of which 12,269 are
coda segments and 9,370 are onset segments. For each segment, the spectral center
of gravity (CoG) is measured along with duration and root-mean-square energy.

\begin{table}[t]
\centering
\small
\setlength{\tabcolsep}{4pt}
\begin{tabular}{llrrrr}
\toprule
Position & Measurement & North & South & $d$ & \makecell{$p$\\(Holm)} \\
\midrule
Coda  & CoG (Hz) & 3,118.7 & 2,417.1 & $+0.56$ & $<$0.001 \\
Coda  & Duration & --      & --      & $+0.30$ & 0.019 \\
Coda  & RMS      & --      & --      & $+0.07$ & 1.000 \\
Onset & CoG (Hz) & 3,517.0 & 3,450.1 & $+0.06$ & 1.000 \\
Onset & Duration & --      & --      & $-0.27$ & 0.033 \\
Onset & RMS      & --      & --      & $+0.03$ & 1.000 \\
\bottomrule
\end{tabular}
\caption{/s/ in North and South, per syllabic position, in speaker level. The
dialectal contrast appears in the coda and disappears in the onset.}
\label{tab:coda}
\end{table}

The pattern is the one predicted by the hypothesis, not the one that would
produce a general difference in sibilant production. In coda, the Southern CoG is
700 Hz lower than the Northern one, with $d = +0.56$. At onset, the difference
between the same varieties is 67 Hz, with $d = +0.06$ and $p = 1.000$. The
interaction between dialect and position on the CoG is significant
($+196.9$ Hz, $p = 0.006$): the South does not produce distinct sibilants;
rather, it produces distinct codas, which is the definition of a phonological
process conditioned by position.

The contrast can be summarized in a single measure per speaker, the weakening
index, defined as the difference between the CoG of their onsets and that of
their codas: the higher the value, the more the coda is weakened relative to the
full realization by the same speaker. The three varieties are ranked as follows:
North, $+359$ Hz over 369 speakers; Mexico, $+836$ Hz over 249; and South,
$+1,022$ Hz over 165.

When comparing the South versus the North and Mexico combined, $d = +0.63$ with
$p < 0.0001$: the strongest effect in the entire study, and the only finding that
dialectology would have predicted in advance. The ranking also replicates what
has been described in the literature without having been imposed by the study
design: Northern peninsular Spanish is the most conservative variety in coda
position, Southern peninsular Spanish is the most innovative, and Mexican Spanish
occupies an intermediate position, consistent with a distribution in which
coastal varieties aspirate and highland varieties maintain
\citep{hualde2005,arias2007}.

\begin{figure*}[t]
\centering
\includegraphics[width=\textwidth]{fig_codas.png}
\caption{Weakening of /s/ in coda: Spectral center of gravity and duration of /s/
according to syllabic position in the three varieties, and F1 of the preceding
vowel according to context. The distinction between varieties appears only in the
coda.}
\label{fig:codas}
\end{figure*}

This is the positive result of the study; its methodological reading is taken up
in Section~\ref{sec:discussion}.

\subsection{The vowel before coda /s/}
\label{sec:precoda}

The compensatory accounts introduced in Section~\ref{sec:dialectal} can now be
tested directly, by contrasting within each speaker the vowels that precede a
coda /s/ with the same vowels in other contexts. A total of 815 speakers provided
data in both contexts: 397 from the North, 168 from the South, and 250 from
Mexico.

The interaction between dialect and context is significant for F1 ($+0.101$,
$p = 0.005$), for F2 ($+0.070$, $p = 0.036$), and for duration ($+0.003$ s,
$p = 0.036$). There is, therefore, a vowel effect associated with the /s/ coda
context that differs across varieties.

Its interpretation, however, requires precision, because it does not match what
is described in the literature. The intra-speaker shift in F1 when moving to a
/s/ coda context is $-0.146$ in the North and $-0.086$ in the South, in Lobanov
z-units. Both values are negative: both varieties close the vowel before the /s/
coda, and the South simply closes it less. There is no Southern vowel opening in
these data; there is a generalized closing in which the South participates to a
lesser extent. The effect is also small, $d = -0.17$, and does not reach
significance on its own ($p = 0.063$). Nor is compensatory lengthening observed:
the difference in duration goes in the same direction but with a smaller
magnitude ($d = -0.10$, $p = 0.241$).

It is worth noting how this result changed with sample size. With 34 Southern
speakers, the same contrast yielded $d = -0.83$, a large effect that fit well
with the traditional description. With 168 speakers, it is $-0.17$. The
interaction remains significant, but the magnitude has been reduced to one-fifth.
This is a textbook case of regression to the mean: small samples do not produce
biased effects in a specific direction; rather, they produce extreme effects, and
those that survive the filter of attention are systematically the largest. It is
the same phenomenon that Section~\ref{sec:pseudo} documents from another angle,
and it is worth keeping this in mind when reading any phonetic study based on
just a few dozen speakers.

\subsection{Exploratory analyses}

Two unsupervised analyses round out the picture and, incidentally, serve to
correct certain conclusions from the preliminary analysis.

By applying $k$-means as implemented in scikit-learn \citep{pedregosa2011} to the
normalized formants and evaluating the resulting partition against the actual
vowel labels using the adjusted Rand index \citep{hubert1985}, normalization
improves the recovery of the vowel categories (ARI = 0.424 with Lobanov versus
0.332 in raw hertz) while the inclusion of bandwidths degrades it (0.357 and
0.258, respectively). Bandwidths do not contribute to vowel structure, they
dilute it.

The preliminary analysis had concluded, based on the silhouette coefficient, that
Spanish is organized into ``three perceptual zones'' rather than five vowels. The
$k$-scan reproduces that result, since the silhouette does peak at $k = 3$
(0.454), but the ARI places it at $k = 4$ (0.479), and the actual number of
categories is five. The discrepancy has a technical, rather than phonological, explanation:
the silhouette measure favors partitions with few compact clusters and decreases
mechanically as dimensions are added, so it is not comparable across spaces of
different dimensionality; the ARI, being measured against a known truth, is
comparable \citep{rousseeuw1987}. The $k = 3$ result was a property of the metric,
not of Spanish.

When training four classifiers to predict dialect based on acoustic parameters,
using speaker-disjoint cross-validation folds, the balanced accuracy is 0.500 in
all four cases, that is, exactly random. Contrasts that include Mexico are only
slightly higher (0.527 and 0.530).

That this reflects the problem and not the characteristics is demonstrated by the
positive control: the same models, applied to the same data and using the same
partitions, predict the vowel identity of never-before-seen speakers with an
accuracy of 0.762 compared to a chance accuracy of 0.200, that is, 3.8 times
chance. The measures contain abundant phonetic information, but what they do not
contain is dialectal information.

The experiment also provides a second demonstration of the argument in
Section~\ref{sec:pseudo}. If the partitions are made randomly across tokens,
allowing the same speaker to appear in both the training and test sets, the
accuracy in dialect prediction is inflated by $+0.20$ on a sample constructed
like the one in the preliminary analysis. On the current sample, the same
comparison yields an increase of $+0.004$. Speaker identity leakage only distorts
the results when a few speakers dominate the sample. When the design is balanced,
both validation schemes converge.

\begin{figure*}[t]
\centering
\includegraphics[width=\textwidth]{fig_classification.png}
\caption{Dialect prediction against vowel identity prediction with
speaker-disjoint cross-validation and with random partitions. Dialect sits at
chance under both schemes in the balanced sample. Vowel identity is recovered far
above chance for previously unseen speakers.}
\label{fig:classification}
\end{figure*}

% =============================================================================
\section{Discussion}
\label{sec:discussion}

In the vowel space, Northern and Southern Peninsular Spanish are
indistinguishable at speaker level. In the treatment of coda /s/ they separate
clearly, with the onset serving as a null control. The large effects that the
preliminary analysis had located in the vowels did not exist at all.

The fact that one of the three results is positive and the other two are null is
no coincidence, and the reason lies in the design rather than in dialectology.
The contrast that survives is the only one that occurs within the speaker.
Comparing the /s/ in coda with the /s/ at the onset of the same speaker instantly
neutralizes everything in this corpus that is confounded with dialect: the
anatomy of the vocal tracts, the microphone, room noise, age, gender, and reading
style. Comparisons between groups, on the other hand, must account for all these
factors through normalization, statistical control, and balanced sampling. The
practical lesson is that, in a corpus of volunteers, a well-chosen intra-speaker
contrast is more valuable than any statistical refinement applied to an overall
comparison.

This also delimits what Common Voice is good for in phonetics. The unevenness
described in Section~\ref{sec:corpus} is not a defect to be tolerated but the
single fact that governs how a study on this corpus must be designed.
Consequently, the speaker identifier is not merely incidental metadata but the
variable that determines whether an analysis is valid. Any study that samples
clips without grouping them by speaker runs the risk of measuring individuals
while believing it is measuring varieties, and this risk increases as the corpus
grows larger, because the volume of data increases statistical significance
without increasing independence. With per-speaker caps and prior aggregation, the
same corpus supports robust inference.

One caveat remains. Section~\ref{sec:vowelspace} established no difference is
detectable, which is not the same as establishing that none exists, and the
difference between those two statements matters here. The descriptive literature
does not describe the Southern vowel system as slightly different from the
Northern one; it describes it as audibly and systematically different. That
stronger claim is the one these data fail to support, in the vowel space as a
whole and also in the compensatory opening before coda /s/, where the effect is
real but five times smaller than a small sample had suggested. The claim they do
support is the consonantal one, and it carries the largest effect in the study.

% =============================================================================
\section{Limitations and conclusions}
\label{sec:limitations}

The data consists of read speech produced by volunteers under uncontrolled
recording conditions; therefore, the conclusions do not automatically apply to
spontaneous speech. Dialect labels are self-reported and unverifiable, and their
granularity is coarse: the label ``Southern Peninsula'' encompasses Andalusian,
Extremaduran, and Murcian varieties whose phonetic behavior is not homogeneous,
which necessarily attenuates any contrast. The sex imbalance is mitigated through
normalization but does not disappear.

In terms of measurement, validation with synthetic vowels reveals a bias of
+33.7\% in the F1 estimated for /u/, so no conclusions should be based on that
vowel; the bias is identical in both dialects and cannot generate spurious
differences, but it can introduce noise. In the analysis of /s/, it is important
to bear in mind that forced alignment sets a boundary for the segment even if it
has been completely elided: a total elision is therefore measured as a segment
with a very low center of gravity rather than as an absence, which likely
underestimates rather than exaggerates the Southern weakening. Finally, the null
result for vowel space constrains but does not eliminate: differences below
$d \approx 0.24$ remain compatible with these data.

Three conclusions, in the order the study reached them.

First, the large effects in the initial analysis (F1 at $d = 1.34$ and B2 at
$d = -1.43$) came from the sampling rather than from the dialects. Two
speakers supplied 54\% and 80\% of their respective cells, and the size of the
effects tracks how much those two are allowed to contribute. Once the speaker
becomes the unit of inference, nothing is left of them.

Second, across 461 Northern and 186 Southern speakers I find no detectable
difference in the vowel space, in a design that would have caught anything at
$d \geq 0.24$. The Mexican outgroup rules out the obvious alternative
explanation: the same procedure, applied to a contrast nobody disputes, produces
effects 2.5 times larger.

Third, what separates these two varieties is consonantal. Southern speakers
weaken coda /s/ substantially and consistently ($d = +0.63$), onset /s/ shows
nothing, and the three varieties line up as the dialectological literature would
predict. This was also the only contrast I tested within speakers rather than
between groups, and I doubt that is a coincidence. On a corpus like this one, a
comparison that each speaker controls for themselves is worth more than a more
elaborate analysis of one that nobody controls.

% =============================================================================
\section*{Declaration on the use of LLMs}

A large language model was used as a writing aid for this report, restricted to
language editing: grammar, spelling, and phrasing of text I had already written,
together with the correction of errors introduced when I translated my own draft
from Spanish into English. It was also used to flag repeated passages so that I
could shorten them.

The research question, the study design, the sampling and analysis code, the
statistical decisions, the results and their interpretation are my own work. No
text was generated from a prompt and inserted without my having written the
corresponding content first, and every number reported here was produced by my
own pipeline and checked against its output tables.

I keep a full record of the editing sessions and can provide it on request.

% =============================================================================
\bibliography{references}
\bibliographystyle{acl_natbib}

\end{document}
